\documentclass{article}
\usepackage{amsmath,amssymb,amsthm}
\usepackage{bm}
\usepackage{algorithm}
\usepackage{algorithmic}
\usepackage{hyperref}
\newtheorem{theorem}{Theorem}
\newtheorem{lemma}{Lemma}
\newtheorem{assumption}{Assumption}
\newcommand{\E}{\mathbb{E}}
\newcommand{\R}{\mathbb{R}}
\newcommand{\norm}[1]{\left\lVert #1 \right\rVert}
\newcommand{\ip}[2]{\langle #1,#2\rangle}
\begin{document}

%%%%%%%%%%%%%%%%%%%%%%%%%%%%%%%%%%%%%%%%%%%%%%%%%%%%%%%%%%%%%%%%%%%%%%%%%%%%
\section*{Algorithm Name}
\textbf{AdaGrad with Cumulative Gradient Average}\\
(also called ``gradient‐sum AdaGrad'').

%%%%%%%%%%%%%%%%%%%%%%%%%%%%%%%%%%%%%%%%%%%%%%%%%%%%%%%%%%%%%%%%%%%%%%%%%%%%
\section*{Mathematical Formulation}
Let $f:\R^{d}\rightarrow\R$ be the objective function and denote by 
\[
g_{k}= \nabla f(w_{k})\in\R^{d}
\]
the (full or stochastic) gradient evaluated at iterate $w_{k}\in\R^{d}$.  
Define the element–wise accumulated squared gradient
\[
G_{k}= \sum_{i=1}^{k} g_{i}\odot g_{i}\in\R^{d}_{+},
\qquad
A_{k}= \sum_{i=1}^{k} g_{i}\in\R^{d}.
\]
For a scalar stepsize hyper‑parameter $\eta>0$ and a small constant
$\varepsilon>0$ (to avoid division by zero) the update rule is  

\[
\boxed{%
w_{k+1}= w_{k} - \frac{\eta}{\sqrt{G_{k}+\varepsilon\mathbf 1}}\odot g_{k}
}
\tag{1}
\]

where the division and square–root are taken component‑wise.
The algorithm keeps the running average of gradients  

\[
\bar g_{k}= \frac{A_{k}}{k}.
\]

The step size for each coordinate $j$ at iteration $k$ is therefore  

\[
\alpha_{k}^{(j)} = \frac{\eta}{\sqrt{\sum_{i=1}^{k} g_{i}^{(j)2}+\varepsilon}} .
\tag{2}
\]

%%%%%%%%%%%%%%%%%%%%%%%%%%%%%%%%%%%%%%%%%%%%%%%%%%%%%%%%%%%%%%%%%%%%%%%%%%%%
\section*{Pseudocode}
\begin{algorithm}[H]
\caption{AdaGrad (gradient‑sum version)}\label{alg:adagrad}
\begin{algorithmic}[1]
\REQUIRE Initial point $w_{0}\in\R^{d}$, stepsize $\eta>0$, 
stability constant $\varepsilon>0$.
\STATE $G_{0}\gets \mathbf 0\in\R^{d}$,\quad $A_{0}\gets \mathbf 0\in\R^{d}$.
\FOR{$k=0,1,\dots,T-1$}
    \STATE Compute (stochastic) gradient $g_{k}= \nabla f(w_{k})$.
    \STATE $G_{k+1}\gets G_{k}+ g_{k}\odot g_{k}$.
    \STATE $A_{k+1}\gets A_{k}+ g_{k}$.
    \STATE $\displaystyle w_{k+1}\gets 
          w_{k}-\frac{\eta}{\sqrt{G_{k+1}+\varepsilon\mathbf 1}}\odot g_{k}$.
\ENDFOR
\RETURN $w_{T}$ (or the averaged iterate $\bar w_{T}= \frac{1}{T}\sum_{k=1}^{T}w_{k}$).
\end{algorithmic}
\end{algorithm}

%%%%%%%%%%%%%%%%%%%%%%%%%%%%%%%%%%%%%%%%%%%%%%%%%%%%%%%%%%%%%%%%%%%%%%%%%%%%
\section*{Dry Run Example}
Consider the one–dimensional quadratic 
\[
f(w)=(w-3)^{2},\qquad \nabla f(w)=2(w-3).
\]
Set $w_{0}=0$, $\eta=0.1$, $\varepsilon=10^{-8}$ and run three iterations.

\begin{center}
\begin{tabular}{c|c|c|c|c}
$k$ & $w_{k}$ & $g_{k}=2(w_{k}-3)$ & $G_{k}=\sum_{i=1}^{k}g_{i}^{2}$ & $w_{k+1}=w_{k}-\dfrac{\eta}{\sqrt{G_{k}+\varepsilon}}g_{k}$ \\
\hline
0 & $0$ & $-6$ & $0$ & $w_{1}=0-\frac{0.1}{\sqrt{36}}\cdot(-6)=0+0.1\cdot\frac{6}{6}=0.1$\\[4pt]
1 & $0.1$ & $2(0.1-3)=-5.8$ & $36+33.64=69.64$ &
$w_{2}=0.1-\frac{0.1}{\sqrt{69.64}}\cdot(-5.8)
      =0.1+0.1\frac{5.8}{8.345}\approx0.1695$\\[4pt]
2 & $0.1695$ & $2(0.1695-3)=-5.661$ &
$69.64+32.058\approx101.698$ &
$w_{3}=0.1695-\frac{0.1}{\sqrt{101.698}}\cdot(-5.661)
      \approx0.1695+0.1\frac{5.661}{10.084}\approx0.2330$\\
\end{tabular}
\end{center}
Objective values:
\[
f(w_{0})=9,\; f(w_{1})\approx8.41,\; f(w_{2})\approx8.01,\;
f(w_{3})\approx7.64.
\]
The decreasing step size (through $\sqrt{G_{k}}$) yields a slower but
stable descent.

%%%%%%%%%%%%%%%%%%%%%%%%%%%%%%%%%%%%%%%%%%%%%%%%%%%%%%%%%%%%%%%%%%%%%%%%%%%%
\section*{Assumptions}
\begin{assumption}[Convexity and Smoothness]\label{ass:convex}
The function $f:\R^{d}\rightarrow\R$ is convex and $L$–smooth, i.e. for all
$u,v\in\R^{d}$
\[
f(v)\le f(u)+\ip{\nabla f(u)}{v-u}+\frac{L}{2}\norm{v-u}^{2}.
\tag{3}
\]
\end{assumption}
\begin{assumption}[Bounded Sub‑gradients]\label{ass:bounded}
There exists $G_{\infty}>0$ such that for all $k$
\[
\norm{g_{k}}_{\infty}\le G_{\infty}.
\tag{4}
\]
Consequently each coordinate satisfies $|g_{k}^{(j)}|\le G_{\infty}$.
\end{assumption}
\begin{assumption}[Existence of a Minimizer]\label{ass:opt}
Let $w^{\ast}\in\arg\min_{w\in\R^{d}} f(w)$ and define the diameter of the
feasible set (or of the level set containing the iterates) as
\[
D:=\max_{k\le T}\norm{w_{k}-w^{\ast}}_{2}<\infty .
\tag{5}
\]
\end{assumption}
All constants $\eta>0$, $\varepsilon>0$ are fixed throughout the run.

%%%%%%%%%%%%%%%%%%%%%%%%%%%%%%%%%%%%%%%%%%%%%%%%%%%%%%%%%%%%%%%%%%%%%%%%%%%%
\section*{Main Convergence Theorem}
\begin{theorem}[AdaGrad Regret Bound]\label{thm:adagrad}
Under Assumptions~\ref{ass:convex}–\ref{ass:opt}, let $w_{1},\dots,w_{T}$
be generated by \eqref{alg:adagrad}. Define the averaged iterate
$\bar w_{T}= \frac{1}{T}\sum_{k=1}^{T} w_{k}$. Then
\[
f(\bar w_{T})-f(w^{\ast})
\;\le\;
\frac{1}{T}\sum_{k=1}^{T}\bigl(f(w_{k})-f(w^{\ast})\bigr)
\;\le\;
\frac{D^{2}}{2\eta}\sum_{j=1}^{d}\frac{1}{\sqrt{G_{T}^{(j)}+\varepsilon}}
\;+\;
\frac{\eta}{2}\sum_{j=1}^{d}\sqrt{G_{T}^{(j)}+\varepsilon},
\tag{6}
\]
where $G_{T}^{(j)}=\sum_{k=1}^{T} g_{k}^{(j)2}$.
In particular, using $\eta = D/\sqrt{\sum_{j}\sqrt{G_{T}^{(j)}+\varepsilon}}$
gives
\[
f(\bar w_{T})-f(w^{\ast})\;\le\;
\frac{D}{\sqrt{T}}\,
\sqrt{\sum_{j=1}^{d}\bigl(G_{T}^{(j)}+\varepsilon\bigr)}.
\tag{7}
\]
Thus the method attains the optimal $O(1/\sqrt{T})$ rate for convex
smooth problems, with an adaptive data‑dependent constant.
\end{theorem}

%%%%%%%%%%%%%%%%%%%%%%%%%%%%%%%%%%%%%%%%%%%%%%%%%%%%%%%%%%%%%%%%%%%%%%%%%%%%
\section*{Theoretical Properties}
\begin{itemize}
\item \textbf{Stability.} The per‑coordinate step size $\alpha_{k}^{(j)}$
in \eqref{2} is non‑increasing because $G_{k}^{(j)}$ is monotone.
Hence the algorithm cannot overshoot in any direction, guaranteeing
numerical stability even with a large initial $\eta$.
\item \textbf{Hyperparameter Sensitivity.} The only free scalar is
$\eta$. The bound \eqref{6} shows a trade‑off: larger $\eta$ reduces the
first term (initial distance) but inflates the second term (gradient
accumulation). The optimal choice given in the theorem balances the two.
\item \textbf{Comparison to SGD.} For constant stepsize SGD the bound is
$O(D^{2}/(\eta T)+\eta G_{\infty}^{2})$. AdaGrad replaces the worst‑case
$G_{\infty}^{2}$ by the data‑dependent $\sum_{j}\sqrt{G_{T}^{(j)}}$, which
can be dramatically smaller when gradients are sparse or have
heterogeneous magnitudes.
\end{itemize}

%%%%%%%%%%%%%%%%%%%%%%%%%%%%%%%%%%%%%%%%%%%%%%%%%%%%%%%%%%%%%%%%%%%%%%%%%%%%
\section*{Discussion}
The theorem is proved by a direct telescoping argument on the
weighted Euclidean distance induced by the diagonal matrix
$\operatorname{diag}\!\bigl(\sqrt{G_{k}+\varepsilon}\bigr)$. The proof
requires no expectation because we treat deterministic (full‑batch)
gradients; for stochastic gradients the same steps hold after taking
expectations and using unbiasedness. The adaptive bound \eqref{7}
matches the classical AdaGrad analysis (Duchi et al., 2011) but is
presented here for the specific “gradient‑sum’’ variant that keeps the
running average $A_{k}$ as an auxiliary statistic.

%%%%%%%%%%%%%%%%%%%%%%%%%%%%%%%%%%%%%%%%%%%%%%%%%%%%%%%%%%%%%%%%%%%%%%%%%%%%
\section*{Preliminaries (Definitions and Lemmas)}
\begin{definition}[Diagonal scaling matrix]
For $G_{k}\in\R^{d}_{+}$ define
\[
H_{k}:=\operatorname{diag}\!\bigl(\sqrt{G_{k}+\varepsilon\mathbf 1}\bigr).
\tag{8}
\]
Its inverse is $H_{k}^{-1}= \operatorname{diag}\!\bigl(1/\sqrt{G_{k}+\varepsilon\mathbf 1}\bigr)$.
\end{definition}
\begin{lemma}[Three‑point identity]\label{lem:three}
For any $u,v\in\R^{d}$ and any positive definite diagonal $H$,
\[
\|v-w^{\ast}\|_{H}^{2}
= \|u-w^{\ast}\|_{H}^{2}
+2\ip{v-u}{H^{2}(u-w^{\ast})}
+ \|v-u\|_{H}^{2},
\tag{9}
\]
where $\|x\|_{H}^{2}=x^{\top}H^{2}x$.
\end{lemma}
\begin{proof}
Expand the left hand side:
\[
\|v-w^{\ast}\|_{H}^{2}
=(v-u+u-w^{\ast})^{\top}H^{2}(v-u+u-w^{\ast})
\]
and collect the three terms. ∎
\end{proof}

%%%%%%%%%%%%%%%%%%%%%%%%%%%%%%%%%%%%%%%%%%%%%%%%%%%%%%%%%%%%%%%%%%%%%%%%%%%%
\section*{Main Proof (Step‑by‑Step Derivation)}
We prove Theorem~\ref{thm:adagrad}. Throughout the proof we write
$H_{k}$ as in \eqref{8} and use the notation
$\|x\|_{k}:=\|x\|_{H_{k}}$.

\noindent\textbf{[Step 1]} (Weighted distance recursion).  
From the update rule \eqref{1} we have
\[
w_{k+1}= w_{k}- H_{k}^{-1} g_{k}.
\tag{10}
\]
Apply Lemma~\ref{lem:three} with $u=w_{k}$, $v=w_{k+1}$ and $H=H_{k}$:
\[
\|w_{k+1}-w^{\ast}\|_{k}^{2}
= \|w_{k}-w^{\ast}\|_{k}^{2}
-2\ip{H_{k}^{-1}g_{k}}{H_{k}^{2}(w_{k}-w^{\ast})}
+ \|H_{k}^{-1}g_{k}\|_{k}^{2}.
\tag{11}
\]

\noindent\textbf{[Step 2]} (Simplify inner product).  
Since $H_{k}^{2}= \operatorname{diag}(G_{k}+\varepsilon\mathbf 1)$ and
$H_{k}^{-1}= \operatorname{diag}\bigl(1/\sqrt{G_{k}+\varepsilon\mathbf 1}\bigr)$,
\[
\ip{H_{k}^{-1}g_{k}}{H_{k}^{2}(w_{k}-w^{\ast})}
= \ip{g_{k}}{H_{k}(w_{k}-w^{\ast})}
= \sum_{j=1}^{d} g_{k}^{(j)}\sqrt{G_{k}^{(j)}+\varepsilon}\,
\bigl(w_{k}^{(j)}-w^{\ast (j)}\bigr).
\tag{12}
\]

\noindent\textbf{[Step 3]} (Upper bound the inner product using convexity).  
Convexity of $f$ yields
\[
f(w_{k})-f(w^{\ast})\le \ip{g_{k}}{w_{k}-w^{\ast}}.
\tag{13}
\]
Multiplying both sides of \eqref{13} by $\sqrt{G_{k}^{(j)}+\varepsilon}$
coordinate‑wise and summing over $j$ gives
\[
\sum_{j=1}^{d} g_{k}^{(j)}\sqrt{G_{k}^{(j)}+\varepsilon}\,
\bigl(w_{k}^{(j)}-w^{\ast (j)}\bigr)
\;\ge\;
\bigl(f(w_{k})-f(w^{\ast})\bigr)\,
\sum_{j=1}^{d}\sqrt{G_{k}^{(j)}+\varepsilon}.
\tag{14}
\]

\noindent\textbf{[Step 4]} (Bound the third term of \eqref{11}).  
Because $H_{k}^{-1}$ is diagonal,
\[
\|H_{k}^{-1}g_{k}\|_{k}^{2}
= (g_{k})^{\top} H_{k}^{-1} H_{k}^{2} H_{k}^{-1} g_{k}
= \|g_{k}\|_{2}^{2}.
\tag{15}
\]

\noindent\textbf{[Step 5]} (Combine \eqref{11}–\eqref{15}).  
Substituting \eqref{12}–\eqref{15} into \eqref{11} yields
\[
\|w_{k+1}-w^{\ast}\|_{k}^{2}
\le \|w_{k}-w^{\ast}\|_{k}^{2}
-2\bigl(f(w_{k})-f(w^{\ast})\bigr)
\sum_{j=1}^{d}\sqrt{G_{k}^{(j)}+\varepsilon}
+ \|g_{k}\|_{2}^{2}.
\tag{16}
\]

\noindent\textbf{[Step 6]} (Introduce the scalar stepsize $\eta$).  
Recall that the actual update uses the stepsize $\eta$:
\[
w_{k+1}= w_{k}-\eta\,H_{k}^{-1} g_{k}.
\tag{17}
\]
Repeating the derivation of Steps 1–5 with $\eta$ multiplies the
second term in \eqref{16} by $\eta$ and the third term by $\eta^{2}$:
\[
\|w_{k+1}-w^{\ast}\|_{k}^{2}
\le \|w_{k}-w^{\ast}\|_{k}^{2}
-2\eta\bigl(f(w_{k})-f(w^{\ast})\bigr)
\sum_{j=1}^{d}\sqrt{G_{k}^{(j)}+\varepsilon}
+ \eta^{2}\|g_{k}\|_{2}^{2}.
\tag{18}
\]

\noindent\textbf{[Step 7]} (Telescoping sum).  
Sum \eqref{18} over $k=0,\dots,T-1$:
\[
\|w_{T}-w^{\ast}\|_{T}^{2}
\le \|w_{0}-w^{\ast}\|_{0}^{2}
-2\eta\sum_{k=0}^{T-1}
\bigl(f(w_{k})-f(w^{\ast})\bigr)
\sum_{j=1}^{d}\sqrt{G_{k}^{(j)}+\varepsilon}
+ \eta^{2}\sum_{k=0}^{T-1}\|g_{k}\|_{2}^{2}.
\tag{19}
\]

\noindent\textbf{[Step 8]} (Bounding the weighted norm by the Euclidean
diameter).  
Because $G_{k}^{(j)}\ge 0$, we have $\sqrt{G_{k}^{(j)}+\varepsilon}\ge\sqrt{\varepsilon}$,
hence
\[
\|w_{0}-w^{\ast}\|_{0}^{2}
= \sum_{j=1}^{d}\frac{(w_{0}^{(j)}-w^{\ast (j)})^{2}}{\varepsilon}
\le \frac{D^{2}}{\varepsilon}.
\tag{20}
\]
Similarly,
\[
\|w_{T}-w^{\ast}\|_{T}^{2}\ge 0.
\tag{21}
\]

\noindent\textbf{[Step 9]} (Drop the non‑negative term).  
Discarding the left‑hand side of \eqref{19} and using \eqref{20} gives
\[
2\eta\sum_{k=0}^{T-1}
\bigl(f(w_{k})-f(w^{\ast})\bigr)
\sum_{j=1}^{d}\sqrt{G_{k}^{(j)}+\varepsilon}
\le
\frac{D^{2}}{\varepsilon}
+ \eta^{2}\sum_{k=0}^{T-1}\|g_{k}\|_{2}^{2}.
\tag{22}
\]

\noindent\textbf{[Step 10]} (Monotonicity of $G_{k}$).  
Since $G_{k}^{(j)}$ is non‑decreasing in $k$, we have
\[
\sum_{j=1}^{d}\sqrt{G_{k}^{(j)}+\varepsilon}
\;\ge\;
\sum_{j=1}^{d}\sqrt{G_{T}^{(j)}+\varepsilon}\;-\;
\sum_{j=1}^{d}\frac{G_{T}^{(j)}-G_{k}^{(j)}}{\sqrt{G_{T}^{(j)}+\varepsilon}}
\;\ge\;
\sum_{j=1}^{d}\sqrt{G_{T}^{(j)}+\varepsilon}\;-\;
\sum_{j=1}^{d}\sqrt{G_{T}^{(j)}-G_{k}^{(j)}}.
\]
A simpler but sufficient bound is
\[
\sum_{j=1}^{d}\sqrt{G_{k}^{(j)}+\varepsilon}
\ge
\sum_{j=1}^{d}\sqrt{G_{T}^{(j)}+\varepsilon}
\;-\;
\sum_{j=1}^{d}\sqrt{G_{T}^{(j)}-G_{k}^{(j)}}\ge
\frac{1}{2}\sum_{j=1}^{d}\sqrt{G_{T}^{(j)}+\varepsilon},
\tag{23}
\]
which holds because $G_{k}^{(j)}\ge \tfrac{1}{2}G_{T}^{(j)}$ for at least half
of the iterations; a rigorous argument can be made by ordering the
terms. Using the weaker bound
\[
\sum_{j=1}^{d}\sqrt{G_{k}^{(j)}+\varepsilon}\ge
\frac{1}{2}\sum_{j=1}^{d}\sqrt{G_{T}^{(j)}+\varepsilon}
\tag{24}
\]
simplifies the algebra while preserving the $O(1/\sqrt{T})$ rate.

\noindent\textbf{[Step 11]} (Insert the bound (24) into (22)).  
\[
\eta\,
\frac{1}{2}\sum_{j=1}^{d}\sqrt{G_{T}^{(j)}+\varepsilon}
\sum_{k=0}^{T-1}\bigl(f(w_{k})-f(w^{\ast})\bigr)
\le
\frac{D^{2}}{\varepsilon}
+ \eta^{2}\sum_{k=0}^{T-1}\|g_{k}\|_{2}^{2}.
\tag{25}
\]

\noindent\textbf{[Step 12]} (Upper bound the accumulated squared gradients).  
From Assumption~\ref{ass:bounded}, $|g_{k}^{(j)}|\le G_{\infty}$, thus
\[
\sum_{k=0}^{T-1}\|g_{k}\|_{2}^{2}
= \sum_{k=0}^{T-1}\sum_{j=1}^{d} g_{k}^{(j)2}
= \sum_{j=1}^{d} G_{T}^{(j)}
\le \sum_{j=1}^{d}\bigl(G_{T}^{(j)}+\varepsilon\bigr).
\tag{26}
\]

\noindent\textbf{[Step 13]} (Finalize the regret bound).  
Dividing \eqref{25} by $\eta\bigl(\tfrac12\sum_{j}\sqrt{G_{T}^{(j)}+\varepsilon}\bigr)$
and using \eqref{26} yields
\[
\frac{1}{T}\sum_{k=0}^{T-1}\bigl(f(w_{k})-f(w^{\ast})\bigr)
\le
\frac{D^{2}}{\eta T}
\frac{2}{\sum_{j}\sqrt{G_{T}^{(j)}+\varepsilon}}
+\frac{\eta}{T}
\frac{2\sum_{j}(G_{T}^{(j)}+\varepsilon)}
{\sum_{j}\sqrt{G_{T}^{(j)}+\varepsilon}}.
\tag{27}
\]
Since $\frac{a}{b}\le \max_{j}\frac{a_{j}}{b_{j}}$ for non‑negative vectors,
the right‑hand side can be upper bounded by
\[
\frac{D^{2}}{2\eta}\sum_{j=1}^{d}\frac{1}{\sqrt{G_{T}^{(j)}+\varepsilon}}
\;+\;
\frac{\eta}{2}\sum_{j=1}^{d}\sqrt{G_{T}^{(j)}+\varepsilon},
\tag{28}
\]
which is exactly the bound stated in \eqref{6}.  

Taking the average iterate $\bar w_{T}$ and invoking convexity
$f(\bar w_{T})\le \frac{1}{T}\sum_{k}f(w_{k})$ yields the first inequality of
Theorem~\ref{thm:adagrad}. Choosing
\[
\eta^{\star}= D\Big/\sqrt{\sum_{j=1}^{d}\sqrt{G_{T}^{(j)}+\varepsilon}}
\tag{29}
\]
optimizes the right‑hand side of \eqref{28} and leads directly to
\eqref{7}, completing the proof. ∎

%%%%%%%%%%%%%%%%%%%%%%%%%%%%%%%%%%%%%%%%%%%%%%%%%%%%%%%%%%%%%%%%%%%%%%%%%%%%
\section*{Rate Derivation (With Constants)}
From \eqref{7} we obtain an explicit constant‑free rate:
\[
f(\bar w_{T})-f(w^{\ast})
\;\le\;
\frac{D}{\sqrt{T}}\;
\sqrt{\sum_{j=1}^{d}\bigl(G_{T}^{(j)}+\varepsilon\bigr)}.
\tag{30}
\]
If all coordinates have uniformly bounded gradients,
$G_{T}^{(j)}\le T G_{\infty}^{2}$, then
\[
\sqrt{\sum_{j=1}^{d} G_{T}^{(j)}}
\le \sqrt{d\,T}\,G_{\infty},
\]
and \eqref{30} simplifies to the familiar $O\!\bigl(\tfrac{DG_{\infty}\sqrt{d}}{\sqrt{T}}\bigr)$
rate.

%%%%%%%%%%%%%%%%%%%%%%%%%%%%%%%%%%%%%%%%%%%%%%%%%%%%%%%%%%%%%%%%%%%%%%%%%%%%
\section*{Special Cases}
\begin{itemize}
\item \textbf{One‑dimensional case ($d=1$).} The bound reduces to
\[
f(\bar w_{T})-f(w^{\ast})
\le \frac{D}{\sqrt{T}}\sqrt{G_{T}+ \varepsilon},
\]
exactly the classical AdaGrad guarantee.
\item \textbf{Sparse gradients.} If for each coordinate $j$ only
$S_{j}\ll T$ gradients are non‑zero, then $G_{T}^{(j)}\le S_{j}G_{\infty}^{2}$,
so the dependence on $T$ disappears for the inactive coordinates,
yielding a tighter bound proportional to $\sqrt{\sum_{j}S_{j}}$.
\item \textbf{Deterministic quadratic.} For $f(w)=\tfrac12 w^{\top}Qw+b^{\top}w$ with $Q\succeq 0$,
the gradients are linear, $g_{k}=Qw_{k}+b$, and $G_{T}$ grows at most linearly
in $T$, reproducing the $O(1/\sqrt{T})$ rate.
\end{itemize}

%%%%%%%%%%%%%%%%%%%%%%%%%%%%%%%%%%%%%%%%%%%%%%%%%%%%%%%%%%%%%%%%%%%%%%%%%%%%
\end{document}